% !TEX root = ../main.tex
\section{Введение}

В настоящей письменной работе применены следующие термины с соответствующими определениями:

Покер --- карточная игра, в которой участники делают ставки в зависимости от силы своих карт и предполагаемых действий соперников~[1].

Онлайн-покер --- форма покера, реализуемая с использованием программных клиентов в сети Интернет~[1].

Раздача --- завершённый игровой цикл, включающий распределение карт, раунды торговли и определение победителя~[1].

Префлоп --- стадия игры, на которой игроки принимают решения до открытия общих карт~[1].

Диапазон рук --- множество стартовых комбинаций карт, с которыми игрок выполняет определённое действие в заданной ситуации~[2].

Позиция за столом --- положение игрока относительно дилера, определяющее порядок хода и влияющее на стратегию игры~[1].

Блайнды --- обязательные ставки, которые делают игроки перед началом раздачи~[1].

Фолд --- действие, при котором игрок отказывается от продолжения участия в раздаче~[1].

Колл --- действие, при котором игрок уравнивает текущую ставку~[1].

Рейз --- действие, при котором игрок увеличивает текущую ставку~[1].

3-бет --- повторное повышение ставки после предыдущего рейза~[3].

RFI --- действие открытия торговли рейзом при отсутствии предыдущих ставок~[1].

HUD --- графический интерфейс, отображаемый поверх основного приложения и предоставляющий дополнительную информацию пользователю~[4].

OCR --- технология распознавания текста и символов на изображениях~[5].

Онлайн-покер является сложной интеллектуальной игрой, сочетающей элементы теории вероятностей, психологии и стратегического планирования. В процессе игры участник принимает решения в условиях неполной информации, ограниченного времени и активного противодействия со стороны соперников.

Одним из ключевых элементов современной стратегии является использование префлоп-диапазонов, представляющих собой заранее определённые наборы стартовых рук для различных игровых ситуаций. Общее количество возможных стартовых комбинаций достигает 169 типов, что формирует значительную когнитивную нагрузку на игрока.

В практической игре это приводит к ряду проблем: сложности запоминания диапазонов, повышенной вероятности ошибок при быстром принятии решений и отклонениям от оптимальной стратегии под влиянием человеческого фактора. Данные проблемы особенно актуальны для начинающих игроков.

Целью данной работы является разработка программного обеспечения, предназначенного для помощи игроку в принятии решений на стадии префлопа в онлайн-покере.

Разрабатываемая система должна обеспечивать анализ текущего состояния игрового стола на основе изображения экрана, определение ключевых параметров игровой ситуации, сопоставление их с заранее заданными диапазонами рук и формирование рекомендации по оптимальному действию.

Система должна функционировать в режиме, близком к реальному времени, обеспечивая минимальную задержку обработки данных, а также поддерживать пользовательские диапазоны и обеспечивать удобный пользовательский интерфейс.

В основе работы используются методы компьютерного зрения и обработки данных. Задача распознавания элементов интерфейса формализуется как задача детекции объектов на изображении и решается с применением свёрточных нейронных сетей. Задача выбора действия сводится к сопоставлению текущего состояния системы с множеством заранее заданных стратегий.

Существующие методы решения включают статические чарты диапазонов, обучающие тренажёры и HUD-системы. Однако данные подходы либо не интегрированы в реальный игровой процесс, либо не предоставляют прямых рекомендаций по действиям в режиме реального времени.

Разрабатываемое программное обеспечение направлено на объединение анализа игровой ситуации и стратегических рекомендаций непосредственно в процессе игры.
